\akapit Zgodnie z instrukcją wykonano następujące czynności:
\begin{enumerate}[itemsep = 1mm, parsep = 1mm, topsep = 3mm, partopsep = 0mm, leftmargin = 1.4cm]
    \item Wybrano parametry pracy tunelu aerodynamicznego i ustawiono zakres kątów dla których miał się odbyć pomiar
    \item Uruchomiono program pomiarów
    \item Wyniki sczytano na pendrive 
\end{enumerate}

\akapit Następnie dokonano niezbędnych obliczeń w celu otrzymania pożądanych charakterystyk płata w funkcji kąta natarcia. Wyniki przedstawiono w sekcji \ref{sec:Wyniki}. Obliczeń dokonano z użyciem środowiska Matlab\footnote{i GNU Octave, ale to kolega ma zboczenie takie}.